\section{新对撞机数据的影响}

\label{sec:impactnewdata}

本节研究 NNPDF3.1 PDF 集合对其所依据的实验信息的依赖。首先我们把新数据的影响与方法学变化的影响区分开。然后我们系统量化 NNPDF3.1 新纳入的每一条实验信息对 PDF 的影响。最后讨论基于特定数据子集的 PDF 确定：
仅用对撞机数据的 PDF（即排除全部固定靶数据）、
仅用质子数据的 PDF（即排除全部核数据），以及排除全部 LHC 数据的 PDF。
由于这些基于缩减数据集的 PDF 集合对特定唯象应用也可能有用，它们同样被公开提供（见下文第~\ref{sec:delivery} 节）。与上一节一样，这里只展示有代表性的图表；更大规模的图表集可在线获取，见第~\ref{sec:delivery} 节所述。


\subsection{分离新数据与方法学的效应}
\label{sec:disentangling}

第~\ref{sec:results-mc} 节研究了 NNPDF3.1 引入的主要方法学改进——独立参数化粲 PDF 并由数据确定——的影响。
为彻底分离数据与方法学的效应，我们用 NNPDF3.1 方法但对 NNPDF3.0 数据集做了一次 PDF 确定：具体做法是从 NNPDF3.1 数据集中去掉所有新数据。该受限数据集与 NNPDF3.0 数据集之间仍有一些小的残余差别，特别是切割上的细微不同，以及使用合并 HERA 数据而非分离的 HERA-I 与 HERA-II 数据。不过这些差异预期很小~\cite{Rojo:2015nxa}。


%%%%%%%%%%%%%%%%%%%%%%%%%%%%%%%%%%%%%%%%%%%%%%%%%%%%%%%%%%%%%%%%%%%%%
\begin{figure}[t]
\begin{center}
  \includegraphics[scale=1]{images/plots/distances_nnpdf31_newdata.pdf}
  \caption{\small 同图~\ref{fig:distances_31_vs_30}，但此处比较的是
    NNPDF3.1 NNLO 全局 PDF 与采用完全相同方法学但以 NNPDF3.0 数据集确定的 PDF。
    \label{fig:distances_nnpdf31_newdata}
  }
\end{center}
\end{figure}
%%%%%%%%%%%%%%%%%%%%%%%%%%%%%%%%%%%%%%%%%%%%%%%%%%%%%%%%%%%%%%%%%%%%%%%

图~\ref{fig:distances_nnpdf31_newdata} 给出 NNPDF3.1 NNLO PDF 集合与同方法学、NNPDF3.0 数据集版本之间的距离。
%
可见新数据的影响主要局限于大 $x$ 区的 up、down、charm 夸克与胶子，
以及中等 $x$ 的奇异味。
%
就不确定度而言，我们在很宽的 $x$ 范围和所有 PDF 味道上都观察到高达半个标准差的改进。
图~\ref{fig:31-nnlo-old-vs-new} 比较了若干代表性 PDF：NNPDF3.1、NNPDF3.0 数据配 NNPDF3.1 方法的集合、以及原始 NNPDF3.0。可以看到新数据与新方法学的总体效应相当，但二者作用于不同区域和不同 PDF。
例如对轻夸克与胶子，在 $x\lsim 10^{-2}$ 区域新方法学的影响占主导，产生增强效应——特别是第~\ref{sec:PDFcomparisons} 节讨论过的 $x\lsim 0.03$ 胶子的增强。
而在大 $x$ 占主导的则是新数据：它导致胶子减小、夸克增强。当然粲受方法学变化影响极大——它在 NNPDF3.0 中未被独立参数化——但对 $x\gsim0.1$，新数据同样有很大影响。
事实上奇异味主要受中等与小 $x$ 区新数据的影响，而粲与胶子在大 $x$ 受其影响最大。

%%%%%%%%%%%%%%%%%%%%%%%%%%%%%%%%%%%%%%%%%%%%%%%%%%%%%%%%%%%%%%%%%%%%%
\begin{figure}[t]
\begin{center}
  \includegraphics[scale=0.38]{images/plots/xg-31-nnlo-old-vs-new.pdf}
  \includegraphics[scale=0.38]{images/plots/xu-31-nnlo-old-vs-new.pdf}
  \includegraphics[scale=0.38]{images/plots/xd-31-nnlo-old-vs-new.pdf}
  \includegraphics[scale=0.38]{images/plots/xdbar-31-nnlo-old-vs-new.pdf}
   \includegraphics[scale=0.38]{images/plots/xsp-31-nnlo-old-vs-new.pdf}
   \includegraphics[scale=0.38]{images/plots/xc-31-nnlo-old-vs-new.pdf}
   \caption{\small 同图~\ref{fig:31-nnlo-vs30}，但另加入以 NNPDF3.1 方法配 NNPDF3.0
     数据确定的 PDF。自左至右、自上而下依次为 gluon、up、down、antidown、total strangeness 与 charm。
    \label{fig:31-nnlo-old-vs-new}
  }
\end{center}
\end{figure}
%%%%%%%%%%%%%%%%%%%%%%%%%%%%%%%%%%%%%%%%%%%%%%%%%%%%%%%%%%%%%%%%%%%%

\subsection{$Z$ 玻色子的横向动量}
\label{sec:impactzpt}
%
长期以来横向动量分布一直被倡导（如文献~\cite{Forte:2013wc}）作为干净而强有力的 PDF 约束，尤其是对胶子。
如第~\ref{sec:expdata} 节所述，得益于该过程直至 NNLO QCD 的计算以及 ATLAS、CMS 8 TeV 的精密 $Z$ $p_T$
数据，如今可以在 NNLO 纳入这类数据。该数据集对 PDF 的影响最近已在文献~\cite{Boughezal:2017nla}中详细研究。

 NNPDF3.1 是首个纳入该数据的全局 PDF 确定。为评估其影响，我们重复了剔除全部 $Z$ $p_T$ 数据的 NNLO 确定。
图~\ref{fig:distances_noZpT} 给出该 PDF 集合与默认集合的距离：显然对所有 PDF 影响温和，变化低于三分之一标准差。
如预期，最大差异出现在胶子与奇异味分布上。造成这种局面的原因最好通过直接比较 PDF 及其不确定度来理解，
见图~\ref{fig:31-nnlo-ZpT}。清楚可见中心值几乎不动而不确定度略有减小，
从而证明这些测量的约束与既有数据集高度一致。因此 $Z$ $p_T$
数据集增强了我们胶子确定的可靠性，也在一定程度上减小了总奇异味的的不确定度。
%
虽然文献~\cite{Boughezal:2017nla}发现该数据集的影响比这里展示的更强，但应指出那是基于 NNPDF3.0 数据集（去除喷注数据）确定 PDF 的情形。NNPDF3.1 增加了更多数据，特别是顶夸克对微分分布：
把 $Z$ $p_T$ 数据加到更宽先验之上时影响变小并不意外。

%%%%%%%%%%%%%%%%%%%%%%%%%%%%%%%%%%%%%%%%%%%%%%%%%%%%%%%%%%%%%%%%%%%%%
\begin{figure}[t]
\begin{center}
  \includegraphics[scale=1]{images/plots/distances_noZpT.pdf}
  \caption{\small 同图~\ref{fig:distances_31_vs_30}，但此处比较默认
    NNPDF3.1 与不包含 ATLAS 和 CMS 8 TeV $Z$ $p_T$ 数据的版本。
\label{fig:distances_noZpT}}
\end{center}
\end{figure}
%%%%%%%%%%%%%%%%%%%%%%%%%%%%%%%%%%%%%%%%%%%%%%%%%%%%%%%%%%%%%%%%%%%%%%%


%%%%%%%%%%%%%%%%%%%%%%%%%%%%%%%%%%%%%%%%%%%%%%%%%%%%%%%%%%%%%%%%%%%%%
\begin{figure}[t]
\begin{center}
  \includegraphics[scale=0.38]{images/plots/xg-31-nnlo-ZpT.pdf}
  \includegraphics[scale=0.38]{images/plots/xsp-31-nnlo-ZpT.pdf}
  \includegraphics[scale=0.38]{images/plots/xg-ERR-31-nnlo-ZpT.pdf}
  \includegraphics[scale=0.38]{images/plots/xsp-ERR-31-nnlo-ZpT.pdf}
  \caption{\small 上排同图~\ref{fig:31-nnlo-vs30}，下排同
    图~\ref{fig:ERR-31-nnlo-vs30}，但此处比较默认 NNPDF3.1 与不含 ATLAS、CMS 8 TeV $Z$ $p_T$ 数据的版本。结果给出 gluon（左）与 total strangeness（right)。
\label{fig:31-nnlo-ZpT}}
\end{center}
\vskip-0.5cm
\end{figure}
%%%%%%%%%%%%%%%%%%%%%%%%%%%%%%%%%%%%%%%%%%%%%%%%%%%%%%%%%%%%%%%%%%%%%%%

除 ATLAS 与 CMS 的 8 TeV 测量外，还存在 ATLAS 7 TeV 的归一化分布测量。纳入该数据集有问题，因为归一化分布的协方差矩阵依赖于对数据集施加的切割，而目前只有完整数据集的协方差矩阵可用。该问题已在文献~\cite{Boughezal:2017nla}中详细研究。此外该数据集已被更精确的 8~TeV 测量取代。因此未将其纳入 NNPDF3.1 数据集。不过，我们通过把它纳入一个专门的 PDF 确定研究了其潜在影响——尽管施加了切割仍沿用其发表的名义协方差矩阵。

表~\ref{tab:NNLOZpTchi2} 给出所有 LHC $Z$ $p_T$ 测量的 $\chi^2/N_{\rm dat}$ 数值，既包括 NNPDF3.1 NNLO 基线（含 ATLAS 与 CMS 8 TeV $Z$ $p_T$ 数据），也包括再纳入 ATLAS 7 TeV $Z$ $p_T$ 数据的确定。
    %
    第一列中 7 TeV 数据的 $\chi^2/N_{\rm dat}$ 置于括号内，以表明与其他数值不同它是预言而非拟合结果。
    显然默认 NNPDF3.1 集合对 ATLAS 7~TeV $Z$ $p_T$ 数据重现很差，即便把它纳入数据集也无法容纳。
    事实上其纳入伴随着对更精确且已取而代之的 ATLAS 8~TeV 数据拟合质量的恶化。
    此外还有迹象表明该数据集与 ATLAS
    $W/Z$ 快度分布之间存在张力：后者的总 $\chi^2$ 恶化了
    $12$ 个单位（46 个数据点）。
这两个 PDF 集合之间的距离示于图~\ref{fig:distances_effect_Zpt7}，
表明纳入 ATLAS 7~TeV $Z$ $p_T$ 数据使胶子与夸克移动了近一倍标准差。
图~\ref{fig:31-nnlo-AZPT7TEV} 对胶子与 down 夸克明确展示了这一点。可以看出不确定度却几乎不变。



%%%%%%%%%%%%%%%%%%%%%%%%%%%%%%%%%%%%%%%%%%%%%%%%%%%%%%%%%%%%%%%%%%%%%%
\begin{table}[b!]
  \centering
  \small
  \begin{tabular}{|c|c|c|}
    \hline
   &         NNPDF3.1 NNLO   &  + ATLAS $Z$ $p_T$ 7 TeV data \\
    \hline
    \hline
ATLAS $Z$ $p_T$ 7 TeV $(p_T^{ll},y_{ll})$    &     [6.78]           &   3.40    \\
ATLAS $Z$ $p_T$ 8 TeV $(p_T^{ll},M_{ll})$    &       0.93         &    0.98   \\
ATLAS $Z$ $p_T$ 8 TeV $(p_T^{ll},y_{ll})$     &    0.93            &   1.17    \\
CMS $Z$ $p_T$ 8 TeV $(p_T^{ll},M_{ll})$     &      1.32          &   1.33    \\
\hline
  \end{tabular}
  \caption{\small
    用 NNPDF3.1 NNLO PDF 集合计算的 LHC $Z$ $p_T$ 数据 $\chi^2/N_{\rm dat}$，
    以及同时纳入 ATLAS $Z$ $p_T$ 7 TeV 数据的新 PDF 确定的相应值。
    \label{tab:NNLOZpTchi2}
  }
\end{table}
%%%%%%%%%%%%%%%%%%%%%%%%%%%%%%%%%%%%%%%%%%%%%%%%%%%%%%%%%%%%%%%%%%%%%%%


虽然无法断言对协方差矩阵更好的处理会如何改变结果，但必须承认：在当前认识水平下，
纳入 ATLAS 7~TeV $Z$ $p_T$ 数据集将对 PDF 产生显著影响而无精度收益，且有迹象显示它与其余 $Z$ $p_T$ 数据集及其他 $W$、$Z$ 产生数据存在张力。因此将其纳入全局数据集并无正当性。

%%%%%%%%%%%%%%%%%%%%%%%%%%%%%%%%%%%%%%%%%%%%%%%%%%%%%%%%%%%%%%%%%%%%%
\begin{figure}[t]
\begin{center}
  \includegraphics[scale=1]{images/plots/distances_effect_Zpt7.pdf}
  \caption{\small 同图~\ref{fig:distances_31_vs_30}，但此处比较默认
    NNPDF3.1 与同时纳入 7~TeV ATLAS $Z$ $p_T$ 数据的版本。
 \label{fig:distances_effect_Zpt7}}
\end{center}
\end{figure}
%%%%%%%%%%%%%%%%%%%%%%%%%%%%%%%%%%%%%%%%%%%%%%%%%%%%%%%%%%%%%%%%%%%%%%%

%%%%%%%%%%%%%%%%%%%%%%%%%%%%%%%%%%%%%%%%%%%%%%%%%%%%%%%%%%%%%%%%%%%%%
\begin{figure}[t]
\begin{center}
  \includegraphics[scale=0.38]{images/plots/xg-31-nnlo-AZPT7TEV.pdf}
  \includegraphics[scale=0.38]{images/plots/xd-31-nnlo-AZPT7TEV.pdf}
  \caption{\small 
同图~\ref{fig:31-nnlo-vs30}，但此处比较默认 NNPDF3.1 与同时纳入
7~TeV ATLAS $Z$ $p_T$ 数据的版本。结果给出 gluon（左）与 down quark（右）。
\label{fig:31-nnlo-AZPT7TEV}}
\end{center}
\end{figure}
%%%%%%%%%%%%%%%%%%%%%%%%%%%%%%%%%%%%%%%%%%%%%%%%%%%%%%%%%%%%%%%%%%%%%%%

\subsection{顶夸克对产生的微分分布}
\label{sec:impacttop}

顶夸克对产生对 PDF 的影响及最优顶夸克数据选择已在文献~\cite{Czakon:2016olj}中广泛讨论。
这里通过比较从数据集中剔除顶夸克数据的 PDF，简要研究顶夸克数据对 NNPDF3.1 的影响。
图~\ref{fig:distances_notop} 给出两集合间的距离。在 $x\gsim 0.1$ 可见胶子中心值与不确定度的巨大差异：
这些数据对大到 $x\simeq 0.6$~\cite{Czakon:2016olj} 的胶子都有约束能力，而其他过程在该区域没有约束力。对其他 PDF 影响温和，最大影响见于小 $x$ 粲。

%%%%%%%%%%%%%%%%%%%%%%%%%%%%%%%%%%%%%%%%%%%%%%%%%%%%%%%%%%%%%%%%%%%%%
\begin{figure}[t]
\begin{center}
  \includegraphics[scale=1]{images/plots/distances_notop.pdf}
  \caption{\small 
同图~\ref{fig:distances_noZpT} 但剔除全部顶夸克数据（总截面与微分布分）。注意左图 $y$ 轴标度不同。
    \label{fig:distances_notop}
  }
\end{center}
\end{figure}
%%%%%%%%%%%%%%%%%%%%%%%%%%%%%%%%%%%%%%%%%%%%%%%%%%%%%%%%%%%%%%%%%%%%%%%
两个 PDF 集合的差异见图~\ref{fig:31-nnlo-top}（图中画出胶子与粲夸克）。大 $x$ 胶子不确定度显著缩小，
无顶数据时的中心值明显高于含顶数据结果的狭窄误差带。这提示顶夸克数据显著提高了胶子确定的精度。
因此图~\ref{fig:31-nnlo-vs30}--\ref{fig:31-nnlo-old-vs-new} 中所见 NNPDF3.1 与 NNPDF3.0 大 $x$ 胶子的差异部分由顶夸克数据驱动。
%
对夸克 PDF 的影响是边缘性的，粲的情形即可见一斑。

正如第~\ref{sec:topdata} 节已提到的，文献~\cite{Czakon:2016vfr}证明快度分布对顶质量的敏感性极小。
事实上文献~\cite{Czakon:2016olj}表明若顶质量变动 1~GeV，
LHC~8~TeV 归一化快度分布的 NLO 理论预言在数据覆盖的运动学范围内最多变化 0.6\%，
远小于数据不确定度或 NNLO 修正的大小。这强烈提示我们的结果本质上与顶质量取值无关。

%%%%%%%%%%%%%%%%%%%%%%%%%%%%%%%%%%%%%%%%%%%%%%%%%%%%%%%%%%%%%%%%%%%%%
\begin{figure}[t]
\begin{center}
  \includegraphics[scale=0.38]{images/plots/xg-31-nnlo-top.pdf}
  \includegraphics[scale=0.38]{images/plots/xc-31-nnlo-top.pdf}
  \includegraphics[scale=0.38]{images/plots/xg-ERR-31-nnlo-top.pdf}
  \includegraphics[scale=0.38]{images/plots/xc-ERR-31-nnlo-top.pdf}
  \caption{\small 
同图~\ref{fig:31-nnlo-ZpT} 但剔除全部顶夸克数据（总截面与微分分布）。
结果给出 gluon（左）与 charm（右），上排为 PDF、下排为其不确定度。
    %
\label{fig:31-nnlo-top}}
\end{center}
\end{figure}
%%%%%%%%%%%%%%%%%%%%%%%%%%%%%%%%%%%%%%%%%%%%%%%%%%%%%%%%%%%%%%%%%%%%%%%

%\clearpage

\subsection{单举喷注产生}
\label{sec:jetdata}

喷注数据用于 PDF 确定由来已久，但其完整的 NNLO 处理直到近期相关计算完成~\cite{Currie:2016bfm,Currie:2017ctp}才成为可能。
然而如第~\ref{sec:datajets} 节所述，NNLO
修正尚未覆盖 NNPDF3.1 纳入的全部数据集。因此在默认 NNPDF3.1 PDF 确定中，喷注采用
NNLO PDF
演化配合 NLO 矩阵元、辅以标度变化确定的额外理论不确定度的方式纳入。这里总体评估喷注数据的效应，特别是这一近似可能的影响。

为此，我们首先重复剔除喷注数据的 NNPDF3.1 确定。所得 PDF 与默认集合的距离示于
图~\ref{fig:distances_nojets}。显然喷注数据的影响温和且非常局域：作用于 $0.1\lsim x\lsim 0.6$ 区域的胶子，
至多半个标准差，对其他 PDF 几乎无影响。其他 PDF 的变化与统计涨落相容。
图~\ref{fig:jetdata-nnlo} 中胶子 PDF 及其不确定度的直接比较证实了这一点。
喷注数据使该区域胶子不确定度最多缩小一半，默认集合的中心值落在较窄的不确定度带内。

有意思的是，在 NNPDF3.0 中喷注数据的影响更为显著：对所有 $x\gsim 0.1$ 不确定度缩小了一个很大的因子。
如第~\ref{sec:impacttop} 节所述，NNPDF3.1 中大 $x$ 胶子受到顶夸克数据的强约束。
具体而言，加入喷注数据在该区域几乎不改变胶子（见图~\ref{fig:jetdata-nnlo}），
而加入顶夸克数据则产生明显位移（见图~\ref{fig:31-nnlo-top}）。
这提示喷注数据与顶夸克数据高度相容，大 $x$ 胶子如今主要由顶夸克数据决定。
这也解释了下文将讨论的对喷注产生 NNLO 修正的不敏感性。


%%%%%%%%%%%%%%%%%%%%%%%%%%%%%%%%%%%%%%%%%%%%%%%%%%%%%%%%%%%%%%%%%%%%%
\begin{figure}[t]
\begin{center}
  \includegraphics[scale=1]{images/plots/distances_nojets.pdf}
  \caption{\small 同图~\ref{fig:distances_noZpT} 但剔除全部喷注数据。
    \label{fig:distances_nojets}
  }
\end{center}
\end{figure}
%%%%%%%%%%%%%%%%%%%%%%%%%%%%%%%%%%%%%%%%%%%%%%%%%%%%%%%%%%%%%%%%%%%%%%%

尽管影响有所减弱，喷注数据仍然扮演着不可忽略的角色。因此有人会担心基于 NLO 矩阵元加理论不确定度的处理是否可靠。
为评估这一点，我们对已有 NNLO 结果的 2011
7~TeV LHC 喷注数据用严格 NNLO 理论重复了 PDF 确定。
其余喷注数据集，即 CDF Run II $k_t$ 喷注、ATLAS 与 CMS 的 $\sqrt{s}=2.76$ TeV 数据集以及 ATLAS 2010
7~TeV 数据集，仍按基线处理。如图~\ref{fig:jetdataex-nnlo}（给出胶子与 down PDF）所示，
PDF 几乎没有任何变化。
该结果与我们以喷注 $p_T$ 为中心标度时发现的百分比水平 NNLO 修正相一致（见
第~\ref{sec:datajets} 节）。

另外，如第~\ref{sec:datajets} 节所述，ATLAS 2011 7~TeV 数据只纳入了中心快度 bin，
因为我们发现虽然逐个纳入各快度 bin 或把不确定度在各 bin 间去关联都能得到良好描述，
但要在保留关联的情况下同时对全部快度 bin 实现良好描述是不可能的。于是自然要问：纳入其他快度 bin 是否会得到不同的 PDF 结果——尽管它们的 PDF 敏感性较低~\cite{Rojo:2014kta}。
为核查这一点，我们用默认 NNPDF3.1 集合把预言与全部 ATLAS 2011 7~TeV 数据比较，并逐个快度 bin 计算 $\chi^2$。
对未纳入的五个快度 bin（从中心到前向），得 $\chi^2/N_{\rm dat}=1.27,\>0.95,\>1.06,\>0.97,\>0.73$，
对应 $N_{\rm dat}=29,\>26,\>23,\>19,\>12$；与之对照，被纳入的中心快度 bin 在表~\ref{tab:NNLOjetschi2}
中的数值为 $N_{\rm dat}=31$ 时 $\chi^2/N_{\rm dat}=1.06$。
我们断定所有快度 bin 都被很好重现，因而任何一个都不可能具有足以在纳入后改变结果的显著拉动。



%%%%%%%%%%%%%%%%%%%%%%%%%%%%%%%%%%%%%%%%%%%%%%%%%%%%%%%%%%%%%%%%%%%%%
\begin{figure}[t]
\begin{center}
  \includegraphics[scale=0.38]{images/plots/xg-31-nnlo-jets.pdf}
  \includegraphics[scale=0.38]{images/plots/xg-ERR-31-nnlo-jets.pdf}
    \caption{\small 默认 NNPDF3.1 NNLO PDF 与剔除全部喷注数据的另一确定之间的比较：
      图中给出 gluon（左）及其百分比不确定度（右）。
\label{fig:jetdata-nnlo}}
\end{center}
\end{figure}

%%%%%%%%%%%%%%%%%%%%%%%%%%%%%%%%%%%%%%%%%%%%%%%%%%%%%%%%%%%%%%%%%%%%%




%%%%%%%%%%%%%%%%%%%%%%%%%%%%%%%%%%%%%%%%%%%%%%%%%%%%%%%%%%%%%%%%%%%%%
\begin{figure}[t]
\begin{center}
  \includegraphics[scale=0.38]{images/plots/xg-31-nnlo-exactjets.pdf}
  \includegraphics[scale=0.38]{images/plots/xd-31-nnlo-exactjets.pdf}
    \caption{\small 
同图~\ref{fig:31-nnlo-vs30}，但此处比较默认 NNPDF3.1 NNLO PDF
      与以严格 NNLO 理论纳入 ATLAS 与 CMS 7~TeV 喷注数据的另一确定。图中给出 gluon（左）与 down（右）PDF。
\label{fig:jetdataex-nnlo}}
\end{center}
\end{figure}

%%%%%%%%%%%%%%%%%%%%%%%%%%%%%%%%%%%%%%%%%%%%%%%%%%%%%%%%%%%%%%%%%%%%%

为更好理解 NNLO 修正及其对 PDF 的影响，
图~\ref{fig:jetdata-nnlo-datath} 将最优拟合预言与 7~TeV 2011 CMS
    与 ATLAS 数据比较，对应
    图~\ref{fig:distances_nojets}--\ref{fig:jetdataex-nnlo} 所比的两个 PDF 集合。
相应的 $\chi^2/N_{\rm dat}$ 收集于表~\ref{tab:NNLOjetschi2}，既包括这些也包括所有其他喷注数据。
首先应注意，对所有尚无严格 NNLO 计算的实验（列于表~\ref{tab:NNLOjetschi2} 上半部分），
两个 PDF 集合给出的 $\chi^2$ 几乎相同。因为这些实验的预言用同一理论计算（NLO 加额外标度不确定度），
这表明 PDF 的变化非常小。
对有 NNLO 理论的实验，$\chi^2$ 数值同样变化甚微。数据与理论的比较（见
图~\ref{fig:jetdata-nnlo-datath}）表明 NLO 与 NNLO 预言如预期相当接近。尽管如此，NNLO
预言与数据符合略好，这反映在表~\ref{tab:NNLOjetschi2} 中更好的 NNLO $\chi^2$ 上——尽管额外标度不确定度（图~\ref{fig:jetdata-nnlo-datath} 中画为数据的内误差棒）
只加在 NLO 预言上。

作为最后的检查，我们重复了带有一个排除大 $p_T$ 喷注数据的切割的 PDF 确定——那里 NNLO 修正更大（对照图~\ref{fig:jetcf}）。
具体地，只保留 $p_T\le 240$~GeV 的 7~TeV 喷注数据、$p_T\le 95$~GeV 的 2.76~TeV
喷注数据、$p_T\le 68$~GeV 的 1.96~TeV
喷注数据。我们发现所得 PDF 与默认集合统计等价（距离量级为一到二）。
我们断定喷注数据近似处理对默认 NNPDF3.1 集合的影响非常小。










%%%%%%%%%%%%%%%%%%%%%%%%%%%%%%%%%%%%%%%%%%%%%%%%%%%%%%%%%%%%%%%%%%%%%%
\begin{table}[H]
  \centering
  \begin{tabular}{|c|c|c|}
    \hline
   &         NNPDF3.1   &  exact NNLO  \\
    \hline
    \hline
CDF Run II $k_t$ jets     &          0.84      &      0.85 \\
ATLAS jets 2.76 TeV   & 1.05       &     1.03 \\
CMS jets 2.76 TeV  & 1.04        &    1.02 \\
ATLAS jets 2010 7 TeV  & 0.96     &       0.95 \\
\hline
ATLAS jets 2011 7 TeV &  1.06      &      0.91 \\
CMS jets 7 TeV 2011 7 TeV &   0.84   &         0.79 \\
\hline
  \end{tabular}
  \caption{\small
    用图~\ref{fig:jetdataex-nnlo} 所比两个 PDF 集合之一得到的全部喷注数据集
    $\chi^2/N_{\rm dat}$，理论取各自 PDF 确定所用者。表中上半部分的数据集尚无严格 NNLO
    计算，全程采用带标度不确定度的 NLO 理论；下半部分右列则用 NNLO 理论。
    \label{tab:NNLOjetschi2}
  }
\end{table}
%%%%%%%%%%%%%%%%%%%%%%%%%%%%%%%%%%%%%%%%%%%%%%%%%%%%%%%%%%%%%%%%%%%%%%

\clearpage


%%%%%%%%%%%%%%%%%%%%%%%%%%%%%%%%%%%%%%%%%%%%%%%%%%%%%%%%%%%%%%%%%%%%%%
\begin{figure}[t]
\begin{center}
  \includegraphics[width=0.49\textwidth]{images/plots/datath-cmsjets.pdf}
  \includegraphics[width=0.49\textwidth]{images/plots/datath-atlasjets.pdf}
  \caption{\small 2011 年 7~TeV CMS（左）与 ATLAS（right)单举单喷注数据与最优拟合结果的比较，
后者分别用带标度不确定度的 NLO 理论或严格 NNLO
理论。最优拟合预言所示不确定度为 PDF 不确定度；数据所示为对角实验误差（外误差棒）
与 NLO 预言标度不确定度（内误差棒）。
\label{fig:jetdata-nnlo-datath}}
\end{center}
\end{figure}
%%%%%%%%%%%%%%%%%%%%%%%%%%%%%%%%%%%%%%%%%%%%%%%%%%%%%%%%%%%%%%%%%%%%%%%


\subsection{前向区的电弱玻色子产生}
\label{sec:lhcbimpact}

%%%%%%%%%%%%%%%%%%%%%%%%%%%%%%%%%%%%%%%%%%%%%%%%%%%%%%%%%%%%%%%%%%%%%
\begin{figure}[t]
\begin{center}
  \includegraphics[scale=1]{images/plots/distances_noLHCb.pdf}
  \caption{\small 同图~\ref{fig:distances_noZpT} 但剔除全部 LHCb 数据。注意左图 $y$ 轴标度不同。
    \label{fig:distances_noLHCb}
  }
\end{center}
\end{figure}
%%%%%%%%%%%%%%%%%%%%%%%%%%%%%%%%%%%%%%%%%%%%%%%%%%%%%%%%%%%%%%%%%%%%%%%

LHCb 实验的电弱产生数据开辟了新的运动学区域，从而为大 $x$ 与小 $x$ 的味分离提供了新的约束。
虽然 LHCb 数据已包含在 NNPDF3.0 中，但涵盖 $W$ 与 $Z$ 数据间全部关联的 Run I 终版 LHCb 7 TeV 与 8 TeV 测量的发布大大提高了其在 PDF 确定中的作用。
图~\ref{fig:distances_noLHCb} 给出默认 NNPDF3.1 NNLO 与剔除全部 LHCb 数据所得 PDF 之间的距离。
%
影响对所有夸克 PDF 都很显著，尤其在价区：
因此这些数据对味分离有实质影响，在大 $x$ 最为突出。



%%%%%%%%%%%%%%%%%%%%%%%%%%%%%%%%%%%%%%%%%%%%%%%%%%%%%%%%%%%%%%%%%%%%%
\begin{figure}[t]
\begin{center}
  \includegraphics[scale=0.34]{images/plots/xu-31-nnlo-LHCb.pdf}
  \includegraphics[scale=0.34]{images/plots/xu-ERR-31-nnlo-LHCb.pdf}
  \includegraphics[scale=0.34]{images/plots/xd-31-nnlo-LHCb.pdf}
  \includegraphics[scale=0.34]{images/plots/xd-ERR-31-nnlo-LHCb.pdf}
  \includegraphics[scale=0.34]{images/plots/xc-31-nnlo-LHCb.pdf}
  \includegraphics[scale=0.34]{images/plots/xc-ERR-31-nnlo-LHCb.pdf}
  \caption{\small 
同图~\ref{fig:31-nnlo-ZpT} 但剔除全部 LHCb 数据。
自上而下给出 up、down 与 charm
PDF。同时显示 PDF（左）与不确定度（right）。
\label{fig:lhcbdata-nnlo}}
\end{center}
\end{figure}
%%%%%%%%%%%%%%%%%%%%%%%%%%%%%%%%%%%%%%%%%%%%%%%%%%%%%%%%%%%%%%%%%%%%%%%
up、down 与 charm PDF 的明确展示见图~\ref{fig:lhcbdata-nnlo}，其中也给出了百分比 PDF
不确定度。LHCb 数据在第~\ref{sec:disentangling} 节讨论过的轻夸克 PDF 数据驱动的大 $x$ 增强
（见图~\ref{fig:31-nnlo-old-vs-new}）中起重要作用，并且是 $x\gsim0.1$ 处新数据对粲产生可观影响的主要原因——在那里它们显著降低了不确定度。
 效应在中等与大 $x$ 更为醒目，峰值在 $x\simeq 0.3$ 附近：该区域 PDF 不确定度也大幅缩小；
 down PDF 的不确定度缩小尤为明显。

为了直接观察 LHCb 数据的影响，
图~\ref{fig:lhcbdata-datatheory} 把 8~TeV LHCb $\mu$ 子 $W^+$ 与 $W^-$ 数据
与用 NNPDF3.0 及 NNPDF3.1 得到的预言比较。改进显而易见，尤其在高快度处。
预言的 PDF 不确定度也有明显缩小。

%%%%%%%%%%%%%%%%%%%%%%%%%%%%%%%%%%%%%%%%%%%%%%%%%%%%%%%%%%%%%%%%%%%%%
\begin{figure}[t]
\begin{center}
  \includegraphics[scale=0.45]{images/plots/LHCb_LHCBWZMU8TEV_dataset_report_datanorm_plot_fancy_0.pdf}
  \includegraphics[scale=0.45]{images/plots/LHCb_LHCBWZMU8TEV_dataset_report_datanorm_plot_fancy_1.pdf}
  \caption{\small 8~TeV LHCb $\mu$ 子 $W^+$（左）与
    $W^-$（right)产生数据和用 NNPDF3.1 及 NNPDF3.0 得到的 NNLO 预言的比较。
    所示不确定度对数据为对角实验不确定度，对最优拟合预言为 PDF 不确定度。
\label{fig:lhcbdata-datatheory}}
\end{center}
\end{figure}
%%%%%%%%%%%%%%%%%%%%%%%%%%%%%%%%%%%%%%%%%%%%%%%%%%%%%%%%%%%%%%%%%%%%%%%

%%%%%%%%%%%%%%%%%%%%%%%%%%%%%%%%%%%%%%%%%%%%%%%%%%%%%%%%%%%%%%%%%%%%%
\begin{figure}[h]
\begin{center}
  \includegraphics[scale=1]{images/plots/distances_d0wasy.pdf}
  \caption{\small
同图~\ref{fig:distances_noZpT} 但剔除 D0 $W$ 不对称度数据。
    \label{fig:distances_nod0wasy}
  }
\end{center}
\end{figure}
%%%%%%%%%%%%%%%%%%%%%%%%%%%%%%%%%%%%%%%%%%%%%%%%%%%%%%%%%%%%%%%%%%%%%%%



\subsection{Tevatron 的 $W$ 不对称度}
\label{sec:legtevdata}

多年来 Tevatron 的 $W$ 产生数据一直是夸克味分解信息的首要来源。
电子道与 $\mu$ 子道的终版 D0 $W$ 不对称度测量已纳入 NNPDF3.1，取代了以往全部数据。
图~\ref{fig:distances_nod0wasy} 对默认 NNPDF3.1 与剔除该数据集的 PDF 进行距离比较。
距离普遍较小，这一观察由图~\ref{fig:d0wasydata-nnlo} 的直接 PDF 比较所印证。
但我们曾在表~\ref{tab:chi2tab_31-nlo-nnlo-30} 看到，该数据集的拟合质量用 NNPDF3.1 比用旧版 NNPDF3.0 好得多。
该数据集影响温和的原因在于它与大量 LHC 数据高度一致——如今正是后者主导着味分离。
因此这些数据进一步佐证了 NNPDF3.1 味分离的可靠性。

%%%%%%%%%%%%%%%%%%%%%%%%%%%%%%%%%%%%%%%%%%%%%%%%%%%%%%%%%%%%%%%%%%%%%
\begin{figure}[h]
\begin{center}
  \includegraphics[scale=0.38]{images/plots/xubar-d0wasydata-nnlo.pdf}
  \includegraphics[scale=0.38]{images/plots/xdbar-d0wasydata-nnlo.pdf}
  \caption{\small 
同图~\ref{fig:31-nnlo-ZpT} 但剔除 D0 $W$
不对称度。
图中给出 antiup（左）与 antidown（right) PDF。
    \label{fig:d0wasydata-nnlo}
  }
\end{center}
\end{figure}
%%%%%%%%%%%%%%%%%%%%%%%%%%%%%%%%%%%%%%%%%%%%%%%%%%%%%%%%%%%%%%%%%%%%

\subsection{ATLAS $W,Z$ 产生数据与奇异味}
\label{sec:atlaswz}

ATLAS $W$ 与 $Z$ 产生数据已包含在 NNPDF3.0 中，但基于
2011 数据集~\cite{Aaboud:2016btc}的近期测量统计不确定度小得多。与先前 ATLAS 测量一样，据称该数据集对奇异味有很大影响。图~\ref{fig:distances_atlaswz11rap} 给出默认 NNPDF3.1 与剔除此数据集版本之间的距离，印证了这一点。
确实，最大的效应——中心值上几乎达一倍标准差——出现在 strange 与 charm PDF 上，
而对所有其他 PDF 影响要小得多。

%%%%%%%%%%%%%%%%%%%%%%%%%%%%%%%%%%%%%%%%%%%%%%%%%%%%%%%%%%%%%%%%%%%%%
\begin{figure}[t]
\begin{center}
  \includegraphics[scale=1]{images/plots/distances_atlaswz11rap.pdf}
  \caption{\small
同图~\ref{fig:distances_noZpT} 但剔除
2011 ATLAS $W,Z$ 快度分布。
    \label{fig:distances_atlaswz11rap}
  }
\end{center}
\end{figure}
%%%%%%%%%%%%%%%%%%%%%%%%%%%%%%%%%%%%%%%%%%%%%%%%%%%%%%%%%%%%%%%%%%%%%%%

图~\ref{fig:pdfs-noAWZrap11} 中 strange 与 charm PDF 的直接比较表明：
纳入 ATLAS 数据使奇异味在中小 $x$ 区显著增强，而粲被压低。如第~\ref{sec:results-mc} 节讨论并如图~\ref{fig:31-nnlo-fitted-vs-pch} 所示，
当粲为微扰生成时这种压低无法容纳；因此参数化粲对于调和 ATLAS 数据与全局数据集至关重要——在该 $x$ 范围全局数据集受到 HERA 数据的强约束。
质子的 strange 与 charm 含量将在下文第~\ref{sec:strangeness} 与 \ref{sec:phenocharm} 节详细讨论。

%%%%%%%%%%%%%%%%%%%%%%%%%%%%%%%%%%%%%%%%%%%%%%%%%%%%%%%%%%%%%%%%%%%%%
\begin{figure}[t]
\begin{center}
  \includegraphics[scale=0.38]{images/plots/xsp-31-nnlo-noAWZrap11.pdf}
  \includegraphics[scale=0.38]{images/plots/xc-31-nnlo-noAWZrap11.pdf}
  \caption{\small 
同图~\ref{fig:31-nnlo-ZpT} 但剔除
2011 ATLAS $W,Z$ 快度分布。
图中给出 total strange（left）与 charm（right) PDF。
    \label{fig:pdfs-noAWZrap11}
  }
\end{center}
\end{figure}
%%%%%%%%%%%%%%%%%%%%%%%%%%%%%%%%%%%%%%%%%%%%%%%%%%%%%%%%%%%%%%%%%%%%


%%%%%%%%%%%%%%%%%%%%%%%%%%%%%%%%%%%%%%%%%%%%%%%%%%%%%%%%%%%%%%%%%%%%%
\begin{figure}[t]
\begin{center}
  \includegraphics[scale=0.42]{images/plots/ATLAS_ATLASWZRAP11_dataset_report_datanorm_plot_fancy_1.pdf}
  \includegraphics[scale=0.42]{images/plots/ATLAS_ATLASWZRAP11_dataset_report_datanorm_plot_fancy_2.pdf}
  \caption{\small 2011 ATLAS 7~TeV $W^-$
    （左）与 $Z$（right)
    数据与用 NNPDF3.1 及 NNPDF3.0 所得 NNLO 预言的比较；$W$
    产生数据按前向轻子赝快度 $\eta_l$ 作图，$Z$ 产生数据按双轻子快度 $y_{ll}$ 作图。
 \label{fig:atlaswz11-datatheory}}
\end{center}
\end{figure}
%%%%%%%%%%%%%%%%%%%%%%%%%%%%%%%%%%%%%%%%%%%%%%%%%%%%%%%%%%%%%%%%%%%%%%%

NNPDF3.1 对数据实现了良好描述，如图~\ref{fig:atlaswz11-datatheory} 所示（图中也给出用 NNPDF3.0 得到的 NNLO 预言）。显然一致性大为改善，
表~\ref{tab:chi2tab_31-nlo-nnlo-30} 的 $\chi^2$ 数值亦证明了这一点。
同样值得注意的是 PDF 不确定度的显著缩小。如第~\ref{data:inclusiveWZ} 节所述，
这些数据只是部分纳入 NNPDF3.1 确定：具体而言偏离峰位或处于前向快度的 $Z$ 产生数据未被纳入。
不过我们核实过，NNPDF3.1 对这些数据的描述同样良好，相对 NNPDF3.0 有类似改善，
$\chi^2/N{\rm dat}$ 量级为一。
未纳入四个 bin 中之二的数据与 NNPDF3.1、NNPDF3.0 的比较示于图~\ref{fig:atlaswz11-datatheorypred}。

%%%%%%%%%%%%%%%%%%%%%%%%%%%%%%%%%%%%%%%%%%%%%%%%%%%%%%%%%%%%%%%%%%%%%
\begin{figure}[t]
\begin{center}
  \includegraphics[scale=0.42]{images/plots/ATLASWZRAP11CC_datanorm_highmasscentral.pdf}
  \includegraphics[scale=0.42]{images/plots/ATLASWZRAP11CF_datanorm_zpeakforward.pdf}
  \caption{\small 同图~\ref{fig:atlaswz11-datatheory}，但对应未纳入
    NNPDF3.1 确定的四个数据 bin 中的两个：中心快度高质量 $Z$ 产生（left）与前向快度壳上 $Z$ 产生（right)。
 \label{fig:atlaswz11-datatheorypred}}
\end{center}
\end{figure}
%%%%%%%%%%%%%%%%%%%%%%%%%%%%%%%%%%%%%%%%%%%%%%%%%%%%%%%%%%%%%%%%%%%%%%%


\subsection{CMS 8~TeV 双微分 Drell-Yan 分布}
\label{sec:cms8tevimpact}

与 ATLAS 类似，CMS 也发表了更新的电弱玻色子产生数据。NNPDF3.0 已纳入来自 CMS 2011 数据集的 7 TeV 双微分（快度与不变质量）Drell-Yan 数据~\cite{Chatrchyan:2013tia}。
%
基于 2012 数据的同一测量在 8~TeV 的更新版发表于文献~\cite{CMS:2014jea}，同时包括绝对截面以及 8~TeV 与 7~TeV 测量之比。

该数据的非关联系统不确定度非常小。遗憾的是只提供了完整协方差矩阵，
而没有分解出各个关联系统误差。这两件事结合在一起使得无法将该实验纳入 NNPDF3.1 数据集，理由如下。
图~\ref{fig:distances_cmsdy2d12} 给出 NNPDF3.1 NNLO PDF 集合与其纳入该数据集的修改版之间的距离。
虽然对不确定度的影响温和，但显然该数据集在对几乎所有 $x$ 的全部 PDF 的中心值层面都有重大影响，
尤以中小 $x$ 胶子为甚。
鉴于 Drell-Yan 产生只能间接触及胶子 PDF，这有些出人意料。

%%%%%%%%%%%%%%%%%%%%%%%%%%%%%%%%%%%%%%%%%%%%%%%%%%%%%%%%%%%%%%%%%%%%%
\begin{figure}[t]
\begin{center}
  \includegraphics[scale=1]{images/plots/distances_cmsdy2d12.pdf}
  \caption{\small
同图~\ref{fig:distances_31_vs_30}，但此处比较默认 NNPDF3.1 与同时纳入
8~TeV CMS 双微分 Drell-Yan 数据的版本。
    \label{fig:distances_cmsdy2d12}
  }
\end{center}
\end{figure}
%%%%%%%%%%%%%%%%%%%%%%%%%%%%%%%%%%%%%%%%%%%%%%%%%%%%%%%%%%%%%%%%%%%%%%%

图~\ref{fig:pdfs-cmsdy2d12} 中 PDF 及其不确定度的直接比较表明：这些数据使 $x\lsim 0.1$ 的胶子上移至多一个标准差，
并使 $x\gsim 0.1$ 的轻夸克 PDF 以相近幅度下移。
然而这并未伴随 PDF 不确定度的缩小——不确定度反而略有增大，如图~\ref{fig:pdfs-cmsdy2d12} 所示。

此外，尽管纳入拟合后 8~TeV CMS
双微分 Drell-Yan 数据自身的拟合质量仍然很差（$\chi^2/N_{\rm dat}=2.88$），
所有其他实验的拟合质量都有一定程度的恶化。
实际上对其余全部数据的总 $\chi^2$ 恶化了 $\Delta \chi^2=11.5$。
更细致的考察显示恶化最明显的是 HERA 合并单举 DIS 数据，$\Delta
\chi^2=19.7$。
这意味着 CMS 数据与全局数据集其余部分之间存在张力，更具体地说与 HERA 数据——即对小 $x$ 胶子最敏感者——存在张力。

我们必须承认该实验似乎与全局数据集不一致，尤其与我们最没有理由怀疑的数据不一致，即合并 HERA
数据及其对胶子的确定。在缺乏关于协方差矩阵更详细信息的情况下无法进一步追究，
因此 8~TeV CMS 双微分 Drell-Yan 数据最终未纳入全局数据集。

%%%%%%%%%%%%%%%%%%%%%%%%%%%%%%%%%%%%%%%%%%%%%%%%%%%%%%%%%%%%%%%%%%%%%
\begin{figure}[t]
\begin{center}
  \includegraphics[scale=0.38]{images/plots/xg-31-nnlo-cmsdy2d12.pdf}
  \includegraphics[scale=0.38]{images/plots/xu-31-nnlo-cmsdy2d12.pdf}
  \includegraphics[scale=0.38]{images/plots/xg-ERR-31-nnlo-cmsdy2d12.pdf}
  \includegraphics[scale=0.38]{images/plots/xu-ERR-31-nnlo-cmsdy2d12.pdf}
  \caption{\small 
上排同图~\ref{fig:31-nnlo-vs30}，但此处比较默认 NNPDF3.1 与同时纳入
8~TeV CMS 双微分 Drell-Yan 数据的版本；下排给出相应百分比不确定度。结果给出 gluon（left）与 up quark（right)。
    \label{fig:pdfs-cmsdy2d12}
  }
\end{center}
\end{figure}
%%%%%%%%%%%%%%%%%%%%%%%%%%%%%%%%%%%%%%%%%%%%%%%%%%%%%%%%%%%%%%%%%%%%

\subsection{EMC $F_2^c$ 数据与内禀粲}
\label{sec:emc}

引入独立参数化粲 PDF 的优点在文献~\cite{Ball:2016neh}中被倡导，那里基于 NNPDF3.0 方法与数据集给出了首个包含粲的全局 PDF 确定。
默认 NNPDF3.0 数据集补充了 EMC 合作组的粲深度非弹性结构函数数据~\cite{Aubert:1982tt}。该数据集相当陈旧，但它仍是大 $x$ 区粲结构函数的唯一测量。
有了更宽的 NNPDF3.1 数据集，EMC 数据集不再那么不可或缺，特别是考虑到 LHC 唯象学：
考虑到人们对其可靠性提出的质疑，它已被排除在默认 NNPDF3.1 确定之外。

不过文献~\cite{Ball:2016neh,Rottoli:2016lsg}中进行的多项检验——诸如运动学切割与系统不确定度的变体——并未显示任何严重的相容性问题，
反而确认该数据集与全局 PDF 确定常规纳入的其他旧固定核靶数据集一样可靠。
因此在 NNPDF3.1 框架下重新审视该数据集的影响是有意义的。
为此我们产生了在默认 NNPDF3.1 数据集上追加 EMC 数据集~\cite{Aubert:1982tt}的全局 NNPDF3.1
NNLO 分析修改版。图~\ref{fig:distances_emc}
给出该 PDF 集合与默认集合之间的距离。
EMC 数据集对粲有不可忽视的一倍标准差量级的影响，并在较小程度（半个标准差量级）上影响所有轻夸克，只有胶子基本不受影响。
效应的主体局限于 $0.01\lsim x\lsim 0.3$ 区域。
PDF 的直接比较见图~\ref{fig:31-nnlo-emc}。
EMC 数据使粲分布在 $0.02\lsim x\lsim 0.2$ 推向默认 PDF 集合中其误差带的上沿，
同时把其不确定度缩小了可观的倍数。轻夸克 PDF 相应地被轻微压低，其不确定度也略有缩小。


%%%%%%%%%%%%%%%%%%%%%%%%%%%%%%%%%%%%%%%%%%%%%%%%%%%%%%%%%%%%%%%%%%%%%
\begin{figure}[t]
\begin{center}
  \includegraphics[scale=1]{images/plots/distances_EMC.pdf}
  \caption{\small 同图~\ref{fig:distances_effect_Zpt7}，但此处比较默认
    NNPDF3.1 与同时纳入 EMC $F_2^c$ 数据集的版本。
    \label{fig:distances_emc}
  }
\end{center}
\end{figure}
%%%%%%%%%%%%%%%%%%%%%%%%%%%%%%%%%%%%%%%%%%%%%%%%%%%%%%%%%%%%%%%%%%%%%%%



表~\ref{tab:EMCchi2} 给出在 NNPDF3.1 数据集中纳入 EMC 数据前后，深度非弹性散射实验及总数据集的 $\chi^2/N_{\rm dat}$。纳入后 EMC 粲数据集的拟合质量大为改善，
其他任何 DIS 数据的拟合质量都没有显著变化：强子对撞机数据敏感性更低。
因此纳入 EMC 粲数据看来能给出更准确的粲确定而在别处无代价；当需要精确的大 $x$ 粲 PDF 时推荐使用该 PDF 集合。粲 PDF 的唯象学意义将在下文第~\ref{sec:phenocharm} 节讨论。

%%%%%%%%%%%%%%%%%%%%%%%%%%%%%%%%%%%%%%%%%%%%%%%%%%%%%%%%%%%%%%%%%%%%%%%
\begin{table}[H]
  \centering
  \small
  \begin{tabular}{|c|c|c|}
    \hline
   &         NNPDF3.1 NNLO   &  + EMC charm data \\
    \hline
    \hline
NMC    &    1.30     &  1.29    \\
SLAC    &    0.75      &  0.76   \\
BCDMS    &    1.21      &  1.24  \\
CHORUS    &    1.11      & 1.10   \\
NuTeV dimuon    &    0.82      &   0.88    \\
\hline
HERA I+II inclusive    &    1.16      &    1.16       \\
HERA $\sigma_c^{\rm NC}$    &    1.45      &  1.42    \\
HERA $F_2^b$    &    1.11      &   1.11   \\
EMC $F_2^c$     &   [4.8]   &   0.93 \\
\hline
\hline
{\bf Total dataset}   &  {\bf  1.148}      &  {\bf  1.145}   \\
\hline
  \end{tabular}
  \caption{\small 深度非弹性散射各实验以及总数据集的 $\chi^2/N_{\rm dat}$，
    分别对应 NNPDF3.1 NNLO PDF 集合与同时纳入 EMC 粲结构函数数据的新 PDF 确定。
    \label{tab:EMCchi2}
  }
\end{table}
%%%%%%%%%%%%%%%%%%%%%%%%%%%%%%%%%%%%%%%%%%%%%%%%%%%%%%%%%%%%%%%%%%%%%%%

\clearpage


%%%%%%%%%%%%%%%%%%%%%%%%%%%%%%%%%%%%%%%%%%%%%%%%%%%%%%%%%%%%%%%%%%%%%
\begin{figure}[t]
\begin{center}
  \includegraphics[scale=0.38]{images/plots/xc-31-nnlo-emc.pdf}
  \includegraphics[scale=0.38]{images/plots/xc-ERR-31-nnlo-emc.pdf}
  \includegraphics[scale=0.38]{images/plots/xu-31-nnlo-emc.pdf}
  \includegraphics[scale=0.38]{images/plots/xd-31-nnlo-emc.pdf}
  \caption{\small 
同图~\ref{fig:31-nnlo-AZPT7TEV}，但此处比较默认
NNPDF3.1 与同时纳入 EMC $F_2^c$ 数据集的版本。
结果给出 charm（左上）、
up（左下）与 down（右下）PDF；charm 的相对 PDF
不确定度亦给出（右上）。
    %
\label{fig:31-nnlo-emc}}
\end{center}
\end{figure}
%%%%%%%%%%%%%%%%%%%%%%%%%%%%%%%%%%%%%%%%%%%%%%%%%%%%%%%%%%%%%%%%%%%%%%%


\subsection{LHC 数据的影响}
\label{sec:nolhc}


我们已经看到 LHC 数据对多种 PDF 有显著影响。既为了精确衡量这种影响，
也因为某些应用需要不带 LHC 数据的 PDF，
我们产生了从 NNPDF3.1 数据集中排除全部 LHC 数据的 PDF 集合。所得集合与默认集合之间的距离示于图~\ref{fig:distances_lhc}。

%%%%%%%%%%%%%%%%%%%%%%%%%%%%%%%%%%%%%%%%%%%%%%%%%%%%%%%%%%%%%%%%%%%%%
\begin{figure}[t]
\begin{center}
  \includegraphics[scale=1]{images/plots/distances_LHC.pdf}
  \caption{\small 同图~\ref{fig:distances_noZpT} 但剔除全部 LHC 数据。
    \label{fig:distances_lhc}
  }
\end{center}
\end{figure}
%%%%%%%%%%%%%%%%%%%%%%%%%%%%%%%%%%%%%%%%%%%%%%%%%%%%%%%%%%%%%%%%%%%%%%%

第~\ref{sec:impactzpt}--\ref{sec:lhcbimpact} 节与第~\ref{sec:atlaswz} 节讨论过的数据的累积效应相当可观。
多数 PDF 受到一倍标准差的影响，个别情形（如 down 与 charm 夸克）甚至达到两倍标准差。
PDF 的直接比较（见图~\ref{fig:31-nnlo-lhc}）印证了这一点。
两次拟合之间的差异似乎主要由 CHORUS、BCDMS 与固定靶 Drell-Yan 数据驱动：剔除 LHC 数据时
它们的 $\chi^2$ 分别改善 84、32 与 38 个单位。其他数据集的差异小得多，
通常与统计涨落相容；某些情形（如 SLAC 数据）全局拟合的拟合质量反而稍好。


另一方面，显然不带 LHC 数据的 PDF 与包含 LHC 数据者之间的位移与各自 PDF 不确定度相容，
且不带 LHC 数据确定的 PDF 其不确定度也没有大到使之对唯象学毫无用处。
我们断定默认集合仍明显更准确，应用于精密唯象学。但在新物理搜寻中，
为确保可能的新物理效应不被重新吸收进 PDF，或许必须使用不带部分或全部 LHC
数据确定的 PDF。
在这类场合，我们的结论是：即便不带 LHC 数据的 PDF 不确定度不具竞争力，其劣化程度也不至于使新物理搜寻完全不可能进行。


%%%%%%%%%%%%%%%%%%%%%%%%%%%%%%%%%%%%%%%%%%%%%%%%%%%%%%%%%%%%%%%%%%%%%
\begin{figure}[t]
\begin{center}
  \includegraphics[scale=0.38]{images/plots/xu-31-nnlo-LHC.pdf}
  \includegraphics[scale=0.38]{images/plots/xd-31-nnlo-LHC.pdf}
  \includegraphics[scale=0.38]{images/plots/xc-31-nnlo-LHC.pdf}
  \includegraphics[scale=0.38]{images/plots/xg-31-nnlo-LHC.pdf}
  \caption{\small 
同图~\ref{fig:31-nnlo-ZpT}
但剔除全部 LHC 数据。结果给出 up（左上）、down（右上）、charm（左下）与
gluon（右下）PDF。
    %
\label{fig:31-nnlo-lhc}}
\end{center}
\end{figure}
%%%%%%%%%%%%%%%%%%%%%%%%%%%%%%%%%%%%%%%%%%%%%%%%%%%%%%%%%%%%%%%%%%%%%%%

\subsection{核靶与核修正}
\label{sec:nonucl}
NNPDF3.1 数据集包含若干在核靶上取得的测量。SLAC、BCDMS 与 NMC 实验的 DIS 数据连同 E886 固定靶 Drell-Yan 数据涉及氘核测量。
全部中微子数据与固定靶 E605 Drell-Yan 数据都用重核靶取得。所有这些数据都已包含在包括 NNPDF3.0 在内的既往 PDF 确定中。
文献~\cite{Ball:2014uwa}研究了核修正的影响，发现可控。
但现在宽泛得多的数据集或许允许把这些数据从全局数据集中移除：移除数据必然带来精度损失，
但这可能因彻底消除对不确定核修正的依赖所带来的准确性提升而得到超额补偿。

为评估这一点，我们用 NNPDF3.1
方法学做了两个额外的 PDF 确定。其一剔除全部重核靶数据但保留氘数据，其二剔除全部核数据只保留质子数据。
默认集合与这两个 PDF 集合之间的距离示于图~\ref{fig:distances_proton}。大 $x$ 处核靶数据的影响显著，在一到两倍标准差水平，
主要体现在海夸克的味分离上。氘数据也有显著影响，特别是在中间 $x$ 范围。

 %%%%%%%%%%%%%%%%%%%%%%%%%%%%%%%%%%%%%%%%%%%%%%%%%%%%%%%%%%%%%%%%%%%%%
\begin{figure}[t]
\begin{center}
  \includegraphics[scale=1]{images/plots/distances_proton1.pdf}\\
   \includegraphics[scale=1]{images/plots/distances_proton2.pdf}
  \caption{\small 同图~\ref{fig:distances_noZpT}，但此处剔除全部重核靶数据而保留
    氘数据（top），或剔除任何核靶数据只保留质子数据（bottom）。
    \label{fig:distances_proton}
  }
\end{center}
\end{figure}
%%%%%%%%%%%%%%%%%%%%%%%%%%%%%%%%%%%%%%%%%%%%%%%%%%%%%%%%%%%%%%%%%%%%%%%

%%%%%%%%%%%%%%%%%%%%%%%%%%%%%%%%%%%%%%%%%%%%%%%%%%%%%%%%%%%%%%%%%%%%%
\begin{figure}[t]
\begin{center}
  \includegraphics[scale=0.35]{images/plots/xg-31-nnlo-proton.pdf}
  \includegraphics[scale=0.35]{images/plots/xu-31-nnlo-proton.pdf}
  \includegraphics[scale=0.35]{images/plots/xd-31-nnlo-proton.pdf}
  \includegraphics[scale=0.35]{images/plots/xdbar-31-nnlo-proton.pdf}
  \caption{\small 
同图~\ref{fig:31-nnlo-ZpT}
但剔除全部重核靶数据保留氘数据，或剔除任何核靶数据只保留质子数据。结果给出 gluon（左上）、up（右上）、down（左下）与 antidown（右下）。
    %
\label{fig:31-nnlo-proton}}
\end{center}
\end{figure}
%%%%%%%%%%%%%%%%%%%%%%%%%%%%%%%%%%%%%%%%%%%%%%%%%%%%%%%%%%%%%%%%%%%%%%%

PDF 的直接比较（图~\ref{fig:31-nnlo-proton}）及其不确定度的比较
（图~\ref{fig:ERR-31-nnlo-datasetvar}）表明：
不带重核靶数据确定的 PDF 与全局集合确实合理相容，只是不确定度偏大，奇异味尤甚。
实际上不带重核靶、甚至不带氘数据的最优拟合结果都在各自不确定度内相互相容，
这与此前结论一致：对这些数据不作核修正在当前 PDF 不确定度水平上不会带来显著偏差。
另一方面只用质子数据确定的 PDF，虽在其更大的不确定度内与全局集合在一倍标准差内相容，
精度损失可观。down 夸克尤为突出，原因在于氘数据对钉扎同位旋三重态 PDF 组合的重要性。

既然氘数据对拟合有显著影响，人们可能担心在当前 PDF 确定的精度下对氘数据的核修正不再可以忽略。
为更细致地考察这一问题，我们做了 NNPDF3.1 NNLO 默认 PDF 确定的一个变体，
其中全部氘数据都加上与 MMHT14 相同的核修正（具体即文献~\cite{Harland-Lang:2014zoa} 式 (9,10)）。



 %%%%%%%%%%%%%%%%%%%%%%%%%%%%%%%%%%%%%%%%%%%%%%%%%%%%%%%%%%%%%%%%%%%%%
\begin{figure}[t]
\begin{center}
  \includegraphics[scale=1]{images/plots/distances_global_deutcorr.pdf}
  \caption{\small 同图~\ref{fig:distances_noZpT}，但此处比较默认
NNPDF3.1 与全部氘数据经文献~\cite{Harland-Lang:2014zoa}核修正的版本。
    \label{fig:distances_nuclcorr}
  }
\end{center}
\end{figure}
%%%%%%%%%%%%%%%%%%%%%%%%%%%%%%%%%%%%%%%%%%%%%%%%%%%%%%%%%%%%%%%%%%%%%%%


%%%%%%%%%%%%%%%%%%%%%%%%%%%%%%%%%%%%%%%%%%%%%%%%%%%%%%%%%%%%%%%%%%%%%
\begin{figure}[t]
\begin{center}
  \includegraphics[scale=0.35]{images/plots/xu-31-nnlo-deuteron.pdf}
  \includegraphics[scale=0.35]{images/plots/xd-31-nnlo-deuteron.pdf}
  \includegraphics[scale=0.35]{images/plots/xu-ERR-31-nnlo-deuteron.pdf}
  \includegraphics[scale=0.35]{images/plots/xd-ERR-31-nnlo-deuteron.pdf}
  \caption{\small 
同图~\ref{fig:31-nnlo-ZpT}
但比较默认 NNPDF3.1 与全部氘数据经文献~\cite{Harland-Lang:2014zoa}核修正的版本。
结果给出 up（左）与 down（right) PDF；不确定度亦给出（bottom row）。
    %
\label{fig:31-nnlo-nuclcorr}}
\end{center}
\end{figure}
%%%%%%%%%%%%%%%%%%%%%%%%%%%%%%%%%%%%%%%%%%%%%%%%%%%%%%%%%%%%%%%%%%%%%%%
就拟合质量而言，我们发现纳入核修正导致轻微恶化：$\chi^2/N_{\rm  dat}=1.156$，
对照默认值 $\chi^2/N_{\rm  dat}=1.148$（见表~\ref{tab:chi2tab_31-nlo-nnlo-30})。
特别地，对 NMC、SLAC 与 BCDMS 数据，加（不加）核修正的 $\chi^2/N_{\rm  dat}$ 分别为 0.94(0.95)、0.71(0.70) 与 1.11(1.11)。
因此添加氘修正对这些数据的拟合质量没有显著影响。

纳入氘修正确定的 PDF 与默认集合的距离示于图~\ref{fig:distances_nuclcorr}。
可以看到影响温和且总低于半个标准差，主要局限于 up 与 down PDF——正如预期。
这些 PDF 示于图~\ref{fig:31-nnlo-nuclcorr}，确认氘修正的影响温和。
应当注意，如图~\ref{fig:31-nnlo-nuclcorr} 所示，纳入氘修正后 PDF 不确定度略有增大。
其他 PDF 的相对位移更小，因为它们本身不确定度较大，而这些不确定度也随核修正的纳入而略有增大。
%（原文此处关于无重核靶版本的对比图被注释，保持原状。）



%%%%%%%%%%%%%%%%%%%%%%%%%%%%%%%%%%%%%%%%%%%%%%%%%%%%%%%%%%%%%%%%%%%%%
%\begin{figure}[t]
%\begin{center}
%  \includegraphics[scale=0.35]{images/plots/xu-31-nnlo-deuteron-nonuclear.pdf}
%  \includegraphics[scale=0.35]{images/plots/xd-31-nnlo-deuteron-nonuclear.pdf}
%  \includegraphics[scale=0.35]{images/plots/xu-ERR-31-nnlo-deuteron-nonuclear.pdf}
%  \includegraphics[scale=0.35]{images/plots/xd-ERR-31-nnlo-deuteron-nonuclear.pdf}
%  \caption{\small 
%Same as Fig.~\ref{fig:31-nnlo-nuclcorr} （原文注释段保持原状）……
%\label{fig:31-nnlo-nuclcorr-nonuc}}
%\end{center}
%\end{figure}
%%%%%%%%%%%%%%%%%%%%%%%%%%%%%%%%%%%%%%%%%%%%%%%%%%%%%%%%%%%%%%%%%%%%%%%

鉴于估计核修正涉及的理论不确定性，并注意到采用文献~\cite{Harland-Lang:2014zoa}模型纳入氘修正后
拟合质量未见改善而 PDF 不确定度略有增加，
我们断定：氘修正对 NNPDF3.1 结果的影响足够小，即使在当前 PDF 确定的高精度下也可以放心地忽略。
尽管如此，未来研究中开展更细致的核修正专门研究——包括与核 PDF 集构造的联系——或许是值得的。

总之，就目前而言在通用 PDF 确定的全局数据集中继续保留核靶数据看起来仍有优势。
但当要求极高精度时（例如标准模型参数的确定），
使用剔除了全部核靶数据的 PDF 集合可能更为妥当。

%%%%%%%%%%%%%%%%%%%%%%%%%%%%%%%%%%%%%%%%%%%%%%%%%%%%%%%%%%%%%%%%%%%%%
\begin{figure}[t]
  \begin{center}
    \includegraphics[scale=0.32]{images/plots/xu-ERR-31-nnlo-datasetvar.pdf}
    \includegraphics[scale=0.32]{images/plots/xubar-ERR-31-nnlo-datasetvar.pdf}
    \includegraphics[scale=0.32]{images/plots/xd-ERR-31-nnlo-datasetvar.pdf}
    \includegraphics[scale=0.32]{images/plots/xdbar-ERR-31-nnlo-datasetvar.pdf}
  \includegraphics[scale=0.32]{images/plots/xs-ERR-31-nnlo-datasetvar.pdf}
  \includegraphics[scale=0.32]{images/plots/xsbar-ERR-31-nnlo-datasetvar.pdf}
  \includegraphics[scale=0.32]{images/plots/xc-ERR-31-nnlo-datasetvar.pdf}
  \includegraphics[scale=0.32]{images/plots/xg-ERR-31-nnlo-datasetvar.pdf}
  \caption{\small $Q=100$ 处 NNPDF3.1 与去重核、仅质子及仅对撞机
    PDF 确定的相对 PDF 不确定度比较。所示不确定度均以
    NNPDF3.1 中心值归一。
    \label{fig:ERR-31-nnlo-datasetvar}
  }
\end{center}
\end{figure}
%%%%%%%%%%%%%%%%%%%%%%%%%%%%%%%%%%%%%%%%%%%%%%%%%%%%%%%%%%%%%%%%%%%%%%%

\subsection{仅对撞机数据的部分子分布}
\label{sec:collideronly}

比上一节更保守的选项是只保留来自 HERA、Tevatron 与 LHC 的对撞机数据。
这一建议最初在 NNPDF2.3 研究~\cite{Ball:2012cx}中提出，动机是这样可排除低标度下取得、
可能受大的微扰与非微扰修正影响的数据。此外还剔除了核靶数据与所有较老的数据集，
从而得到更可靠的 PDF 集合。但由于当时对撞机数据有限，此前的仅对撞机 PDF 具有非常大的不确定度。

为了在当前宽广得多的 LHC 数据集下重新评估形势，我们重复了一次仅对撞机的 PDF 确定。
这相当于重复上一节描述的仅质子确定，再把质子固定靶数据也去掉。
所得 PDF 与默认 NNPDF3.1 的距离示于
图~\ref{fig:distances_collider}。与图~\ref{fig:distances_proton} 对照可见
大多数 PDF 的距离相似，主要例外是胶子——它在中间 $x$ 区现在移动了近两个标准差。

%%%%%%%%%%%%%%%%%%%%%%%%%%%%%%%%%%%%%%%%%%%%%%%%%%%%%%%%%%%%%%%%%%%%%
\begin{figure}[t]
\begin{center}
  \includegraphics[scale=1]{images/plots/distances_nnpdf31_collider.pdf}
  \caption{\small 
同图~\ref{fig:distances_noZpT} 但只保留对撞机数据。
    \label{fig:distances_collider}
  }
\end{center}
\end{figure}
%%%%%%%%%%%%%%%%%%%%%%%%%%%%%%%%%%%%%%%%%%%%%%%%%%%%%%%%%%%%%%%%%%%%%%%

图~\ref{fig:pdfs-collideronly} 的 PDF 直接比较与图~\ref{fig:ERR-31-nnlo-datasetvar}
的不确定度比较证实了这一点。
价夸克（尤其是 up）相当稳定，但在剔除全部固定靶数据后海夸克变得相当不稳定，
明显超过仅质子集合（图~\ref{fig:31-nnlo-proton}）的程度，特别是大 $x$ anti-up 夸克分布的不确定度大幅增加。
此外，在图~\ref{fig:31-nnlo-proton} 的仅质子集合中曾相当稳定的胶子，
现在在中间 $x$ 经历明显的向下移动，尽管其不确定度并未大幅增加。

我们的结论是：尽管近期 LHC 测量带来了令人瞩目的改进，
仅对撞机的 PDF 确定对一般唯象应用而言仍然用处不大。


%%%%%%%%%%%%%%%%%%%%%%%%%%%%%%%%%%%%%%%%%%%%%%%%%%%%%%%%%%%%%%%%%%%%%
\begin{figure}[t]
\begin{center}
  \includegraphics[scale=0.38]{images/plots/xg-31-nnlo-collider.pdf}
  \includegraphics[scale=0.38]{images/plots/xu-31-nnlo-collider.pdf}
  \includegraphics[scale=0.38]{images/plots/xd-31-nnlo-collider.pdf}
  \includegraphics[scale=0.38]{images/plots/xdbar-31-nnlo-collider.pdf}
  \caption{\small 
同图~\ref{fig:31-nnlo-ZpT}
但只保留对撞机数据。结果给出 gluon（左上）、up（右上）、down（左下）与 antidown（右下）。
    \label{fig:pdfs-collideronly}
  }
\end{center}
\end{figure}
%%%%%%%%%%%%%%%%%%%%%%%%%%%%%%%%%%%%%%%%%%%%%%%%%%%%%%%%%%%%%%%%%%%%



